\begin{center}
  \centering
  \captionof{table}{Self-Reference and Limits at a glance: the eighteen laws of \Cref{sec:self-reference-limits}, grouped into reading groups.}
  \label{tab:axis-self-reference-limits-summary}
  \scriptsize
  \setlength{\tabcolsep}{4pt}
  \medskip\noindent\begin{minipage}{\textwidth}\centering
  \textbf{Limits}\par\smallskip
  \begin{tabularx}{\textwidth}{@{}
      >{\raggedright\arraybackslash}p{0.24\textwidth}
      >{\raggedright\arraybackslash}X
      >{\raggedright\arraybackslash}p{0.10\textwidth}@{}}
    \toprule
    Name & One-line claim & Substrate cite \\
    \midrule
    Horizon-Sufficiency Theorem (\Fref{F33}) &
    When the pairing is surjective, a target depends on a history only
    through its boundary readout and
    exterior class exactly when it factors through their pairing; a
    bundled family does so exactly when each member does.  &
    \textemdash \\
    Information-Loss Theorem (\Fref{F34}) &
    When recovery after transport is surjective, source information is
    recoverable after transport exactly when it
    agrees on any two sources that transport and recovery identify.  &
    \textemdash \\
    Singular-Limit Theorem (\Fref{F36}) &
    For a surjective limit quotient, a target is completed exactly when no two processes with the same
    limit object have different target values, and is path-dependent
    otherwise; the pairing \(\rho_{u,\ell}\) refines \(\ell\), and \(u\) is completed
    through it; \(R\circ u\) is completed when \(u\) is, conversely for injective
    \(R\).  &
    \textemdash \\
    Locality Theorem (\Fref{F45}) &
    For a surjective local quotient, a target is local to a domain exactly
    when it is constant on the domain's classes. For strict coarsening, with the signature in the
    declared class and a nonempty carrier, a domain is minimal exactly
    when it has the signature's fibers.  &
    \textemdash \\
    \bottomrule
  \end{tabularx}
  \end{minipage}\par

  \medskip\noindent\begin{minipage}{\textwidth}\centering
  \textbf{Emergence}\par\smallskip
  \begin{tabularx}{\textwidth}{@{}
      >{\raggedright\arraybackslash}p{0.24\textwidth}
      >{\raggedright\arraybackslash}X
      >{\raggedright\arraybackslash}p{0.10\textwidth}@{}}
    \toprule
    Name & One-line claim & Substrate cite \\
    \midrule
    Scale-Descent Theorem (\Fref{F35}) &
    For a surjective scale quotient and coarsening, a law descends to the
    coarser scale exactly when it descends to the
    finer scale and its descended law is constant on the fibers of the
    coarsening; the renormalized law is then unique.  &
    \textemdash \\
    Continuum-Emergence Theorem (\Fref{F43}) &
    Over a nonempty directed index set with a surjective limit quotient,
    a coherent, cofinal tower with vanishing residuals, compatible
    projections and stage readouts, and a target readout equal to one
    stage readout after projection and constant on the limit fibers
    yields \(\operatorname{StableLimitQuotient}(L)\), witnessed by the unique readout on the quotient agreeing with every stage readout; coherence, cofinality and vanishing enter only as conjuncts of StableLimitQuotient, and the content is descent of the common stage readout through \(\ell\).  &
    \textemdash \\
    Fine-Tuning Theorem (\Fref{F47}) &
    If no modulus reaches the target region, fine-tuning is not
    realized; for a monotone measure record with a transitive order, a
    stricter target whose selector region has positive measure keeps a small
    selector region small.  &
    \textemdash \\
    Vacuum Theorem (\Fref{F50}) &
    For a surjective moduli quotient and a descended background, exactly
    one holds: it is necessary
    (equivalently, the raw background is constant) or it varies across
    moduli; a vacuum exists exactly when the vacuum criterion is a
    singleton with formed target, and is then unique.  &
    \textemdash \\
    Unification Theorem (\Fref{F51}) &
    Under exact unification the parent quotient is a common refinement of
    both children; when \(q_U\) is surjective, \(\langle\pi_A,\pi_B\rangle\) is
    injective iff the parent keeps exactly the children's distinctions.  &
    \textemdash \\
    \bottomrule
  \end{tabularx}
  \end{minipage}\par

  \medskip\noindent\begin{minipage}{\textwidth}\centering
  \textbf{Invariants}\par\smallskip
  \begin{tabularx}{\textwidth}{@{}
      >{\raggedright\arraybackslash}p{0.24\textwidth}
      >{\raggedright\arraybackslash}X
      >{\raggedright\arraybackslash}p{0.10\textwidth}@{}}
    \toprule
    Name & One-line claim & Substrate cite \\
    \midrule
    Complementarity Theorem (\Fref{F37}) &
    On a nonempty carrier, two separately admissible accesses are
    complementary exactly when
    every admissible joint carrier identifies two histories that one
    access separates, so that no admissible carrier has both accesses
    factoring through it.  &
    \textemdash \\
    Arrow Theorem (\Fref{F38}) &
    For a surjective record quotient, a continuation has an arrow exactly
    when it descends to a map of record classes that never lowers a
    record's status.  &
    \textemdash \\
    Entropy Theorem (\Fref{F39}) &
    For a surjective access quotient, the target obstruction is empty
    exactly when the target descends. For an entropy functional with a declared zero satisfying
    \(E(q')=0\) exactly when \(q'\) is injective,
    and with the chain rule and combine-monotonicity at a coarsening
    pair, \(E(q)=0\) exactly when \(q\) has no raw split pair, and entropy
    does not decrease under the coarsening. For finite nonempty \(H\), the fiber-volume functional
    \(E(q')=|H|^{-1}\sum_h\log|q'^{-1}(q'h)|\) satisfies these hypotheses and has \(m(\langle q,t\rangle,q)=0\) exactly
    when the target descends.  &
    \textemdash \\
    Irreducibility Theorem (\Fref{F42}) &
    Closedness at \(h\) is emptiness of the target and route split-pair sets on the class of \(h\); with well-founded costs and some admissible closed compression, an irreducible one exists; if for every \(v\in E\) the admissible class contains a map taking \(h\) to \(v\) and isolating \(h\), irreducibility is minimal cost over \(E\) and ignores the target.  &
    \textemdash \\
    Criticality Theorem (\Fref{F44}) &
    A phase change along a finite path of formed controls in which, for every neighbourhood index, each control has a neighbour in the phase of the next, or along a continuous path when neighbourhoods are the open sets, crosses a phase boundary; boundary and criticality
    descend through the control quotient when their data do.  &
    \textemdash \\
    Law--State Split Theorem (\Fref{F46}) &
    For a surjective moduli quotient, a readout descends exactly when no two closures with one modulus differ, and then uniquely; a descended readout is a law readout exactly when it is constant on the closure space and state-dependent exactly when two closures with different moduli differ; the selected modulus fixes the readout on its fiber.  &
    \textemdash \\
    Topology Theorem (\Fref{F48}) &
    For a surjective gluing quotient, a readout factors through it exactly when its split-pair set is empty; each such readout is exactly one of complete and coarse; a transport induces a map of gluing classes exactly when \(\gamma'\circ\tau\) has empty split-pair set.  &
    \textemdash \\
    Common-Source Theorem (\Fref{F49}) &
    On a nonempty history space, common-source and interface
    explanations are equivalent to their factorizations; direct-route
    explanation additionally requires admissibility and stability. A
    correlation residual is exactly failure of local explainability, and a
    common source supplies an interface factorization with \(\iota=\sigma\).  &
    \textemdash \\
    Closure Completeness Theorem (\Fref{F52}) &
    Scoped closure completeness holds exactly when the closure is formed
    and scope-declared, every residual is statused and adequate, and
    every claimed direction is in scope or named outside.  &
    \textemdash \\
    \bottomrule
  \end{tabularx}
  \end{minipage}\par

\end{center}
